\subsection{Decomposition and Complementation}
Throughout this subsection, all vector spaces are over the same fixed field $\mathbb{F}$. Dimension counts first detect directness; complements then realize quotients as concrete subspaces.

\begin{theorem} \label{thm:dimension-formula-for-sum} \hlt{Dimension Formula for a Sum}\\
If $\mathcal{U},\mathcal{W}\leq\mathcal{V}$ are finite-dimensional, then $\dim(\mathcal{U}+\mathcal{W})=\dim(\mathcal{U})+\dim(\mathcal{W})-\dim(\mathcal{U}\cap\mathcal{W})$.
\begin{proof}
The subspace $\mathcal{U}\cap\mathcal{W}$ is finite-dimensional by Proposition \ref{prop:dimension-bounds-and-basis-criteria}. Choose basis $\epsilon=(\mathbf{e}_1,\ldots,\mathbf{e}_r)$ of intersection by Corollary \ref{cor:finite-dimensional-spaces-and-subspaces-have-bases}.\\
Extend $\epsilon$ to bases $\beta_{\mathcal{U}}=(\mathbf{e}_1,\ldots,\mathbf{e}_r,\mathbf{u}_1,\ldots,\mathbf{u}_p)$ of $\mathcal{U}$ and $\beta_{\mathcal{W}}=(\mathbf{e}_1,\ldots,\mathbf{e}_r,\mathbf{w}_1,\ldots,\mathbf{w}_q)$ of $\mathcal{W}$ by Theorem \ref{thm:extending-independent-lists}.\\
Set $\gamma\defeq(\mathbf{e}_1,\ldots,\mathbf{e}_r,\mathbf{u}_1,\ldots,\mathbf{u}_p,\mathbf{w}_1,\ldots,\mathbf{w}_q)$. The list $\gamma$ spans $\mathcal{U}+\mathcal{W}$.\\
Suppose $\sum\limits_{i=1}^{r}a_i\mathbf{e}_i+\sum\limits_{j=1}^{p}b_j\mathbf{u}_j+\sum\limits_{k=1}^{q}c_k\mathbf{w}_k=\mathbf{0}$. Then $\sum\limits_{i=1}^{r}a_i\mathbf{e}_i+\sum\limits_{j=1}^{p}b_j\mathbf{u}_j=-\sum\limits_{k=1}^{q}c_k\mathbf{w}_k\in\mathcal{U}\cap\mathcal{W}$.\\
Its representation in $\beta_{\mathcal{U}}$ therefore has $b_1=\cdots=b_p=0$.\\
The remaining relation belongs to $\beta_{\mathcal{W}}$, so $a_1=\cdots=a_r=c_1=\cdots=c_q=0$.\\
Thus $\gamma$ is a basis of $\mathcal{U}+\mathcal{W}$, and $\dim(\mathcal{U}+\mathcal{W})=r+p+q=(r+p)+(r+q)-r$.\\
\hlt{Structural Commentary}: An intersection basis is extended twice; concatenation then separates shared from new directions.\\
\hlt{Key Insight}: Subtract the dimension counted in both summands.
\end{proof}
\end{theorem}

\begin{corollary} \label{cor:dimension-direct-sum} \hlt{Dimension of a Direct Sum}\\
If $\mathcal{U},\mathcal{W}\leq\mathcal{V}$ are finite-dimensional and $\mathcal{U}\cap\mathcal{W}=\{\mathbf{0}\}$, then $\dim(\mathcal{U}\oplus\mathcal{W})=\dim(\mathcal{U})+\dim(\mathcal{W})$.
\begin{proof}
The sum is direct by Proposition \ref{prop:uniqueness-direct-sum-two-subspaces}. Since $\dim(\mathcal{U}\cap\mathcal{W})=0$, the asserted equality holds by Theorem \ref{thm:dimension-formula-for-sum}.\\
\hlt{Structural Commentary}: The directness criterion removes the intersection term from the sum formula.\\
\hlt{Key Insight}: Trivial overlap makes dimensions additive.
\end{proof}
\end{corollary}

\begin{corollary} \label{cor:large-subspaces-intersect} \hlt{Large Subspaces Must Intersect}\\
Let $\mathcal{U},\mathcal{W}\leq\mathcal{V}$, where $\mathcal{V}$ is finite-dimensional, and let $k\in\N$.
\begin{enumerate}[label=(\roman*)]
\setlength{\itemsep}{0pt}
\item If $\dim(\mathcal{U})+\dim(\mathcal{W})>\dim(\mathcal{V})$, then $\mathcal{U}\cap\mathcal{W}$ contains a nonzero vector.
\item If $\dim(\mathcal{U})+\dim(\mathcal{W})\geq\dim(\mathcal{V})+k$, then $\mathcal{U}\cap\mathcal{W}$ contains $k$ linearly independent vectors.
\end{enumerate}
\begin{proof}
The inequality $\dim(\mathcal{U}+\mathcal{W})\leq\dim(\mathcal{V})$ holds by Proposition \ref{prop:dimension-bounds-and-basis-criteria}.\\
Hence $\dim(\mathcal{U}\cap\mathcal{W})=\dim(\mathcal{U})+\dim(\mathcal{W})-\dim(\mathcal{U}+\mathcal{W})\geq\dim(\mathcal{U})+\dim(\mathcal{W})-\dim(\mathcal{V})$ by Theorem \ref{thm:dimension-formula-for-sum}.
\begin{enumerate}[label=(\roman*)]
\setlength{\itemsep}{0pt}
\item Under strict inequality, $\dim(\mathcal{U}\cap\mathcal{W})\geq1$. Choose basis of intersection by Corollary \ref{cor:finite-dimensional-spaces-and-subspaces-have-bases}; any one of its vectors is nonzero.
\item Under second inequality, $\dim(\mathcal{U}\cap\mathcal{W})\geq k$.\\
A basis of the intersection exists by Corollary \ref{cor:finite-dimensional-spaces-and-subspaces-have-bases}, and any $k$ of its vectors are linearly independent.
\end{enumerate}
\hlt{Structural Commentary}: The sum formula converts excess ambient dimension into a lower bound on overlap.\\
\hlt{Key Insight}: Excess dimension must reappear in the intersection.
\end{proof}
\end{corollary}

\begin{theorem} \label{thm:direct-sum-iff-dimensions-add} \hlt{Direct Sums and Dimension}\\
Let $m\in\N$, $\mathcal{V}$ be finite-dimensional, $\mathcal{U}_1,\ldots,\mathcal{U}_m\leq\mathcal{V}$. For $1\leq j\leq m$, let $\beta_j$ be an ordered basis of $\mathcal{U}_j$, $\beta=(\beta_1,\ldots,\beta_m)$ be their concatenation.
\begin{enumerate}[label=(\roman*)]
\setlength{\itemsep}{0pt}
\item The sum $\mathcal{U}_1+\cdots+\mathcal{U}_m$ is direct iff $\dim(\mathcal{U}_1+\cdots+\mathcal{U}_m)=\sum\limits_{j=1}^{m}\dim(\mathcal{U}_j)$.
\item If the sum is direct, then $\beta$ is an ordered basis of $\mathcal{U}_1\oplus\cdots\oplus\mathcal{U}_m$.
\end{enumerate}
\begin{proof}
The list $\beta$ spans $\mathcal{U}_1+\cdots+\mathcal{U}_m$ and has length $\sum\limits_{j=1}^{m}\dim(\mathcal{U}_j)$.
\begin{enumerate}[label=(\roman*)]
\setlength{\itemsep}{0pt}
\item[$(i)$] ($\Rightarrow$) Suppose the sum is direct. A linear relation among the vectors of $\beta$ determines vectors $\mathbf{u}_j\in\mathcal{U}_j$ satisfying $\sum\limits_{j=1}^{m}\mathbf{u}_j=\mathbf{0}$.\\
Every $\mathbf{u}_j$ is zero by Proposition \ref{prop:internal-direct-sum-family-criterion}, and then every coefficient is zero by the independence of $\beta_j$.\\
Thus $\beta$ is a basis, and its length proves the dimension equality.
\item[$(i)$] ($\Leftarrow$) Suppose dimension equality holds. The spanning list $\beta$ has length $\dim(\mathcal{U}_1+\cdots+\mathcal{U}_m)$, so it is a basis by Proposition \ref{prop:full-dimension-basis-criteria}.\\
If $\sum\limits_{j=1}^{m}\mathbf{u}_j=\mathbf{0}$ with $\mathbf{u}_j\in\mathcal{U}_j$, expanding each $\mathbf{u}_j$ in $\beta_j$ produces a relation in $\beta$.\\
Its coefficients, and hence every $\mathbf{u}_j$, are zero. The sum is therefore direct by Proposition \ref{prop:internal-direct-sum-family-criterion}.
\item[$(ii)$] Under directness, the linear-independence argument in the forward direction of part (i) applies to $\beta$.\\
Since $\beta$ already spans the sum, it is an ordered basis of $\mathcal{U}_1\oplus\cdots\oplus\mathcal{U}_m$.
\end{enumerate}
\hlt{Structural Commentary}: Concatenation converts directness into independence and dimension equality into full basis length.\\
\hlt{Key Insight}: Dimensions add exactly when concatenation has no redundancy.
\end{proof}
\end{theorem}

\begin{definition} \label{def:adapted-basis-direct-sum} \hlt{Basis Adapted to a Direct Sum}\\
Let $m\in\N$, let $\mathcal{V}=\mathcal{U}_1\oplus\cdots\oplus\mathcal{U}_m$ be finite-dimensional, and let $\beta$ be an ordered basis of $\mathcal{V}$.\\
The basis $\beta$ is \hlt{adapted to the direct-sum decomposition} iff there are ordered bases $\beta_j$ of $\mathcal{U}_j$ such that $\beta=(\beta_1,\ldots,\beta_m)$.
\end{definition}

\begin{proposition} \label{prop:complements-exist-finite-dimensions} \hlt{Complements Exist in Finite Dimensions}\\
If $\mathcal{V}$ is finite-dimensional and $\mathcal{U}\leq\mathcal{V}$, then $(\exists\,\mathcal{W}\leq\mathcal{V})\,\mathcal{V}=\mathcal{U}\oplus\mathcal{W}$.
\begin{proof}
Choose a basis $(\mathbf{u}_1,\ldots,\mathbf{u}_m)$ of $\mathcal{U}$ by Corollary \ref{cor:finite-dimensional-spaces-and-subspaces-have-bases}, and extend it to a basis $(\mathbf{u}_1,\ldots,\mathbf{u}_m,\mathbf{w}_1,\ldots,\mathbf{w}_n)$ of $\mathcal{V}$ by Theorem \ref{thm:extending-independent-lists}.\\
Set $\mathcal{W}\defeq\operatorname{span}(\mathbf{w}_1,\ldots,\mathbf{w}_n)$. The extended basis spans $\mathcal{V}$, so $\mathcal{V}=\mathcal{U}+\mathcal{W}$.\\
If $\mathbf{x}\in\mathcal{U}\cap\mathcal{W}$, subtracting its two representations produces a relation in the extended basis.\\
Linear independence forces every coefficient to vanish, so $\mathbf{x}=\mathbf{0}$. Therefore $\mathcal{V}=\mathcal{U}\oplus\mathcal{W}$ by Proposition \ref{prop:uniqueness-direct-sum-two-subspaces}.\\
\hlt{Structural Commentary}: Basis extension isolates the old directions from the newly appended directions.\\
\hlt{Key Insight}: The appended basis vectors span a complement.
\end{proof}
\end{proposition}

\begin{theorem} \label{thm:complements-exist-arbitrary-vector-spaces} \hlt{Complements Exist}\\
Assume Zorn's lemma. If $\mathcal{U}\leq\mathcal{V}$, then $(\exists\,\mathcal{W}\leq\mathcal{V})\,\mathcal{V}=\mathcal{U}\oplus\mathcal{W}$.
\begin{proof}
Choose a Hamel basis $\mathcal{B}_{\mathcal{U}}$ of $\mathcal{U}$ and extend it to a Hamel basis $\mathcal{B}$ of $\mathcal{V}$ by Theorem \ref{thm:hamel-basis-extension}. Set $\mathcal{C}\defeq\mathcal{B}\setminus\mathcal{B}_{\mathcal{U}}$ and $\mathcal{W}\defeq\operatorname{span}(\mathcal{C})$.\\
Every finite $\mathcal{B}$-expansion separates into terms indexed by $\mathcal{B}_{\mathcal{U}}$ and terms indexed by $\mathcal{C}$, so $\mathcal{V}=\mathcal{U}+\mathcal{W}$.\\
If $\mathbf{x}\in\mathcal{U}\cap\mathcal{W}$, then $\mathbf{x}$ has one $\mathcal{B}$-expansion supported in $\mathcal{B}_{\mathcal{U}}$ and another supported in the disjoint set $\mathcal{C}$.\\
These expansions agree by Proposition \ref{prop:hamel-coordinate-representation}, so both have empty support and $\mathbf{x}=\mathbf{0}$. Therefore $\mathcal{V}=\mathcal{U}\oplus\mathcal{W}$ by Proposition \ref{prop:uniqueness-direct-sum-two-subspaces}.\\
\hlt{Structural Commentary}: Hamel basis extension partitions all coordinates into subspace and complementary supports.\\
\hlt{Key Insight}: Disjoint coordinate supports produce a complement.
\end{proof}
\end{theorem}

\begin{proposition} \label{prop:rank-nullity-complement-view} \hlt{Rank-Nullity as a Complement Isomorphism}\\
Let $T\in\LL(\mathcal{V},\mathcal{W})$ and suppose $\mathcal{V}=\ker(T)\oplus\mathcal{M}$.
\begin{enumerate}[label=(\roman*)]
\setlength{\itemsep}{0pt}
\item The restriction $T|_{\mathcal{M}}:\mathcal{M}\to\operatorname{range}(T)$ is an isomorphism.
\item If $\mathcal{V}$ is finite-dimensional, then $\dim(\mathcal{V})=\dim(\ker(T))+\dim(\operatorname{range}(T))$.
\end{enumerate}
\begin{proof}
\begin{enumerate}[label=(\roman*)]
\setlength{\itemsep}{0pt}
\item The restriction is linear and $\ker(T|_{\mathcal{M}})=\mathcal{M}\cap\ker(T)=\{\mathbf{0}\}$, so it is injective by Proposition \ref{prop:injective-iff-kernel-zero}.\\
If $T\mathbf{v}\in\operatorname{range}(T)$, write $\mathbf{v}=\mathbf{k}+\mathbf{m}$ with $\mathbf{k}\in\ker(T)$ and $\mathbf{m}\in\mathcal{M}$. Then $T\mathbf{v}=T\mathbf{m}$, so the restriction is surjective onto $\operatorname{range}(T)$.\\
It is therefore an isomorphism by Proposition \ref{prop:linear-map-invertible-iff-bijective}.
\item Suppose $\mathcal{V}$ is finite-dimensional. The subspaces $\ker(T)$ and $\mathcal{M}$ are finite-dimensional by Proposition \ref{prop:dimension-bounds-and-basis-criteria}, and $\operatorname{range}(T)$ is finite-dimensional by Theorem \ref{thm:range-spanned-by-images-of-basis}. Thus $\dim(\mathcal{V})=\dim(\ker(T))+\dim(\mathcal{M})$ by Corollary \ref{cor:dimension-direct-sum}, while $\dim(\mathcal{M})=\dim(\operatorname{range}(T))$ by part (i) and Theorem \ref{thm:finite-dimensional-isomorphic-iff-same-dimension}. Substitution proves the formula.
\end{enumerate}
\hlt{Structural Commentary}: The restriction removes precisely the invisible kernel directions before dimensions are counted.\\
\hlt{Key Insight}: A kernel complement models the range exactly.
\end{proof}
\end{proposition}

\begin{proposition} \label{prop:finite-dimensional-one-sided-inverses-split} \hlt{One-Sided Inverses Split in Finite Dimensions}\\
Let $\mathcal{V}$ and $\mathcal{W}$ be finite-dimensional, and let $T\in\LL(\mathcal{V},\mathcal{W})$.
\begin{enumerate}[label=(\roman*)]
\setlength{\itemsep}{0pt}
\item If $T$ is injective, then $(\exists\,L\in\LL(\mathcal{W},\mathcal{V}))\,LT=I_{\mathcal{V}}$.
\item If $T$ is surjective, then $(\exists\,R\in\LL(\mathcal{W},\mathcal{V}))\,TR=I_{\mathcal{W}}$.
\end{enumerate}
\begin{proof}
\begin{enumerate}[label=(\roman*)]
\setlength{\itemsep}{0pt}
\item Suppose $T$ is injective. Choose $\mathcal{H}\leq\mathcal{W}$ with $\mathcal{W}=\operatorname{range}(T)\oplus\mathcal{H}$ by Proposition \ref{prop:complements-exist-finite-dimensions}.\\
For the unique decomposition $\mathbf{w}=T\mathbf{v}+\mathbf{h}$, define $L\mathbf{w}\defeq\mathbf{v}$.\\
This is well-defined: directness determines $T\mathbf{v}$ and $\mathbf{h}$, while injectivity determines $\mathbf{v}$.\\
If $\mathbf{w}_i=T\mathbf{v}_i+\mathbf{h}_i$ for $i=1,2$ and $a,b\in\mathbb{F}$, then $L(a\mathbf{w}_1+b\mathbf{w}_2)=L(T(a\mathbf{v}_1+b\mathbf{v}_2)+(a\mathbf{h}_1+b\mathbf{h}_2))=a\mathbf{v}_1+b\mathbf{v}_2=aL\mathbf{w}_1+bL\mathbf{w}_2$.\\
Hence $L$ is linear, and $LT\mathbf{v}=\mathbf{v}$ for every $\mathbf{v}\in\mathcal{V}$.
\item Suppose $T$ is surjective. Choose $\mathcal{H}\leq\mathcal{V}$ with $\mathcal{V}=\ker(T)\oplus\mathcal{H}$ by Proposition \ref{prop:complements-exist-finite-dimensions}.\\
Since $\operatorname{range}(T)=\mathcal{W}$, the restriction $T|_{\mathcal{H}}:\mathcal{H}\to\mathcal{W}$ is an isomorphism by Proposition \ref{prop:rank-nullity-complement-view}(i).\\
Let $R\defeq(T|_{\mathcal{H}})^{-1}$, regarded as a linear map from $\mathcal{W}$ to $\mathcal{V}$. Then $TR=I_{\mathcal{W}}$.
\end{enumerate}
\hlt{Structural Commentary}: A range complement extends a left inverse; a kernel complement supports a right inverse.\\
\hlt{Key Insight}: Complements turn injectivity or surjectivity into a splitting.
\end{proof}
\end{proposition}

\begin{proposition} \label{prop:quotient-by-complement-is-complement} \hlt{Quotient by a Complemented Subspace}\\
Let $\mathcal{U},\mathcal{W}\leq\mathcal{V}$ with $\mathcal{V}=\mathcal{U}\oplus\mathcal{W}$. Then the map $J:\mathcal{W}\to\mathcal{V}/\mathcal{U}$ defined by $J(\mathbf{w})\defeq\mathbf{w}+\mathcal{U}$ is an isomorphism.\\
Consequently, every complement of $\mathcal{U}$ is isomorphic to $\mathcal{V}/\mathcal{U}$.
\begin{proof}
\begin{enumerate}[label=(\roman*)]
\setlength{\itemsep}{0pt}
\item \hlt{Linearity}: For $a,b\in\mathbb{F}$, $\mathbf{w}_1,\mathbf{w}_2\in\mathcal{W}$, equality $J(a\mathbf{w}_1+b\mathbf{w}_2)=aJ(\mathbf{w}_1)+bJ(\mathbf{w}_2)$ holds by Proposition \ref{prop:quotient-space-structure}. Thus $J$ is linear.
\item \hlt{Injectivity}: The equivalence $J(\mathbf{w})=\mathcal{U}\Leftrightarrow\mathbf{w}\in\mathcal{U}$ holds by Proposition \ref{prop:parallel-affine-subsets-equal-or-disjoint}.\\
Hence $\ker(J)=\mathcal{U}\cap\mathcal{W}=\{\mathbf{0}\}$, and $J$ is injective by Proposition \ref{prop:injective-iff-kernel-zero}.
\item \hlt{Surjectivity}: Let $\mathbf{v}+\mathcal{U}\in\mathcal{V}/\mathcal{U}$. Write $\mathbf{v}=\mathbf{u}+\mathbf{w}$ with $\mathbf{u}\in\mathcal{U}$ and $\mathbf{w}\in\mathcal{W}$.\\
Then $\mathbf{v}+\mathcal{U}=\mathbf{w}+\mathcal{U}=J(\mathbf{w})$ by Proposition \ref{prop:parallel-affine-subsets-equal-or-disjoint}, so $J$ is surjective.
\item \hlt{Conclusion}: The map $J$ is an isomorphism by Proposition \ref{prop:linear-map-invertible-iff-bijective}.\\
Since $\mathcal{W}$ was an arbitrary complement of $\mathcal{U}$, this proves the consequent.
\end{enumerate}
\hlt{Structural Commentary}: The proof checks that each coset has exactly one representative in the complement.\\
\hlt{Key Insight}: A complement is a concrete model of the quotient.
\end{proof}
\end{proposition}

\begin{definition} \label{def:codimension-subspace} \hlt{Codimension}\\
If $\mathcal{U}\leq\mathcal{V}$ and $\mathcal{V}/\mathcal{U}$ is finite-dimensional, then the \hlt{codimension of $\mathcal{U}$ in $\mathcal{V}$} is $\operatorname{codim}_{\mathcal{V}}(\mathcal{U})\defeq\dim(\mathcal{V}/\mathcal{U})$.
\end{definition}

If $\mathcal{V}$ is finite-dimensional, then $\operatorname{codim}_{\mathcal{V}}(\mathcal{U})=\dim(\mathcal{V})-\dim(\mathcal{U})$ by Theorem \ref{thm:dimension-of-quotient}. If $\mathcal{V}=\mathcal{U}\oplus\mathcal{W}$ and $\mathcal{V}/\mathcal{U}$ is finite-dimensional, then $\mathcal{W}\cong\mathcal{V}/\mathcal{U}$ by Proposition \ref{prop:quotient-by-complement-is-complement}, so $\operatorname{codim}_{\mathcal{V}}(\mathcal{U})=\dim(\mathcal{W})$ by Proposition \ref{prop:isomorphisms-preserve-bases}.

\begin{perspective} \label{per:quotients-split-over-fields} \hlt{Quotients Split over Fields}\\
In finite dimensions, a complement $\mathcal{V}=\mathcal{U}\oplus\mathcal{W}$ exists by Proposition \ref{prop:complements-exist-finite-dimensions}, and $\mathcal{V}/\mathcal{U}\cong\mathcal{W}$ by Proposition \ref{prop:quotient-by-complement-is-complement}.\\
Thus the quotient forgets the $\mathcal{U}$-component and retains the complementary coordinate.
\end{perspective}

\begin{proposition} \label{prop:quotient-of-direct-sum-by-direct-sum} \hlt{Quotient of a Direct Sum}\\
If $\mathcal{V}=\mathcal{V}_1\oplus\mathcal{V}_2$ and $\mathcal{U}_j\leq\mathcal{V}_j$ for $j=1,2$, then $\mathcal{V}/(\mathcal{U}_1\oplus\mathcal{U}_2)\cong(\mathcal{V}_1/\mathcal{U}_1)\times(\mathcal{V}_2/\mathcal{U}_2)$.
\begin{proof}
\begin{enumerate}[label=(\roman*)]
\setlength{\itemsep}{0pt}
\item \hlt{Denominator}: Since $\mathcal{U}_1\cap\mathcal{U}_2\subseteq\mathcal{V}_1\cap\mathcal{V}_2=\{\mathbf{0}\}$, the sum $\mathcal{U}_1+\mathcal{U}_2$ is direct by Proposition \ref{prop:uniqueness-direct-sum-two-subspaces}.
\item \hlt{Construction}: For the unique decomposition $\mathbf{v}=\mathbf{v}_1+\mathbf{v}_2$, define $\Phi:\mathcal{V}\to(\mathcal{V}_1/\mathcal{U}_1)\times(\mathcal{V}_2/\mathcal{U}_2)$ by $\Phi(\mathbf{v}_1+\mathbf{v}_2)\defeq(\mathbf{v}_1+\mathcal{U}_1,\mathbf{v}_2+\mathcal{U}_2)$.\\
Componentwise quotient operations show that $\Phi$ is linear by Proposition \ref{prop:quotient-space-structure}.
\item \hlt{Range and kernel}: Every pair $(\mathbf{x}_1+\mathcal{U}_1,\mathbf{x}_2+\mathcal{U}_2)$ equals $\Phi(\mathbf{x}_1+\mathbf{x}_2)$, so $\Phi$ is surjective.\\
Moreover, $\Phi(\mathbf{v}_1+\mathbf{v}_2)=(\mathcal{U}_1,\mathcal{U}_2)$ iff $\mathbf{v}_j\in\mathcal{U}_j$ for $j=1,2$ by Proposition \ref{prop:parallel-affine-subsets-equal-or-disjoint}. Hence $\ker(\Phi)=\mathcal{U}_1\oplus\mathcal{U}_2$.
\item \hlt{Conclusion}: The claimed isomorphism holds by Theorem \ref{thm:quotient-by-kernel-is-range}, because $\operatorname{range}(\Phi)$ is the whole product.
\end{enumerate}
\hlt{Structural Commentary}: Direct-sum coordinates define the map; its componentwise kernel identifies the quotient.\\
\hlt{Key Insight}: Quotienting a direct sum acts coordinatewise.
\end{proof}
\end{proposition}

\begin{perspective} \label{per:quotients-and-complements-are-not-functorial-warnings} \hlt{Quotients Depend on the Embedding}\\
A complement realizes a quotient, but the result depends on how the subspace is embedded.\\
Let $\mathcal{V}$ have countable basis $(\mathbf{e}_n)_{n\geq1}$ and set $\mathcal{S}\defeq\operatorname{span}(\mathbf{e}_n:n\geq2)$.\\
The forward shift $F:\mathcal{V}\to\mathcal{S}$ defined by $F(\mathbf{e}_n)=\mathbf{e}_{n+1}$ is an isomorphism: unique finite basis expansions make it linear, and the backward shift $\mathbf{e}_{n+1}\mapsto\mathbf{e}_n$ is its inverse. Thus $\mathcal{S}\cong\mathcal{V}$.\\
Nevertheless, $\mathcal{V}=\mathcal{V}\oplus\{\mathbf{0}\}$ and $\mathcal{V}=\mathcal{S}\oplus\mathbb{F}\mathbf{e}_1$, so $\mathcal{V}/\mathcal{V}\cong\{\mathbf{0}\}$ while $\mathcal{V}/\mathcal{S}\cong\mathbb{F}\mathbf{e}_1$ by Proposition \ref{prop:quotient-by-complement-is-complement}.\\
Hence isomorphic subspaces can have nonisomorphic quotients when embedded differently. 
\end{perspective}
